\subsection{Classical Hopfield Networks}

A classical Hopfield network consists of a collection of interconnected neurons whose connections encode a set of stored patterns. During retrieval, the network gradually updates its internal state until it converges to the stored pattern that most closely resembles the input. In this way, the network functions as an associative memory, where partial or corrupted inputs can trigger the recall of complete stored patterns.

Despite their conceptual importance, classical Hopfield networks have several limitations. Their storage capacity is limited, and retrieval performance deteriorates as more patterns are stored. The original formulation operates on binary representations, making it difficult to represent the complex, high-dimensional data commonly encountered in modern machine learning applications. These limitations motivated the development of modern Hopfield networks, which extend the original framework to continuous representations and substantially larger memory capacities.